\section{Completion of the trace theorem}
\label{sec:completion}

We now combine the geometric vanishing with the detector. The order of
the choices matters: the graph is fixed by the idempotent before its
cycle dimension is chosen.

\begin{proof}[Proof of Theorem~\ref{thm:bass}]
First let $G$ be finitely generated, let $e\in M_q(\C G)$ be an
idempotent, and fix an infinite-order $g\in G$. Construct the graph on
$Y=\langle g\rangle\backslash G$ and the free action of
$H=C_G(g)/\langle g\rangle$ as in Section~\ref{sec:idempotents}.
Let $D\geq2$ be a degree bound for this fixed graph, and choose an even
integer $n=2m\geq n_0(D)$ with $m\geq1$.

Proposition~\ref{prop:path-cycle} gives the cycle $z_n$, and
Proposition~\ref{prop:detector} gives an $H$-invariant cocycle
$\psi_m:E_n\to\C$ with
\[
 \psi_m(z_n)=2^{-m}\lambda,
 \qquad
 \lambda=\HS_G([e(\C G)^q])([g]_G).
\]
If no nonzero simplices occur, $z_n=0$ and this identity already gives
$\lambda=0$. Otherwise Proposition~\ref{prop:path-cycle} realizes $z_n$
as a cycle in $C_n(K)_H$ for an $H$-invariant complex $K$ satisfying all
hypotheses of Theorem~\ref{thm:sparse}. That theorem gives
$\psi_m(z_n)=0$, with the same conclusion. Every finitely generated
projective right $\C G$-module is represented by an idempotent, and every element of
$K_0(\C G)$ is a difference of two such classes. Additivity proves the
assertion for all finitely generated groups.

For an arbitrary discrete group $G$, choose idempotents representing
$x\in K_0(\C G)$ and let $G_0\leq G$ be generated by $g$ and all group
elements occurring in their entries. Then $G_0$ is finitely generated
and $x=i_*x_0$ for some $x_0\in K_0(\C G_0)$, where $i:G_0\hookrightarrow G$.
Lemma~\ref{lem:locality} expresses $\HS_G(x)([g]_G)$ as the sum of the
coefficients $\HS_{G_0}(x_0)(C)$ over those $G_0$-conjugacy classes
$C\subseteq G_0\cap[g]_G$. Only finitely many terms are nonzero, and
every element in every such class has infinite order. Each term
vanishes by the finitely generated case.
\end{proof}

\subsection{The integral trace}

For the integral group ring, the same idempotent formula defines
\[
 \HS_{\Z G}:K_0(\Z G)\longrightarrow
 \Z G/[\Z G,\Z G]\cong\bigoplus_{C\in\con(G)}\Z C,
\]
where commutators are now taken with their additive span. The integral
Bass trace conjecture asks for support only at the identity class.

\begin{corollary}\label{cor:integral}
For every discrete group $G$, the Hattori--Stallings trace of every
class in $K_0(\Z G)$ is supported on $[1]_G$.
\end{corollary}

\begin{proof}
If $P$ is a finitely generated projective right $\Z G$-module, the same
idempotent matrix over $\C G$ represents $P\otimes_{\Z G}\C G$.
Consequently its trace coefficients are the images under
$\Z\hookrightarrow\C$ of those of $P$. Theorem~\ref{thm:bass} and this
coefficientwise injection give vanishing at all infinite-order elements.
No injectivity assertion about the map between the two $K_0$ groups is
required.

For finite-order elements we use Linnell's restriction
\cite[Lemma 4.1]{Linnell1983}, in the form stated and reproved by
Berrick and Hesselholt \cite[Theorem A]{BH2015}: if the integral trace of a finitely
generated projective module is nonzero at $h$, there is an integer
$a\geq1$ such that $h$ is conjugate to $h^{s^a}$ for every integer
$s\geq1$. Their convention also uses right modules.
If $h$ has finite order $t>1$, take $s=t$.
Then $h^{t^a}=1$, so the restriction would make $h$ conjugate to $1$,
a contradiction. Thus every nonidentity finite-order coefficient also
vanishes. Taking differences of projective classes finishes the proof.
\end{proof}

\subsection{Homotopy idempotents and Lefschetz numbers}

The integral support statement also has fixed-point consequences.
A self-map $f$ is a \emph{homotopy idempotent} if $f\circ f\simeq f$.
Such a map on a connected finite CW complex $X$ admits maps
$u:X\to Y$ and $v:Y\to X$ with
\[
 v\circ u\simeq f,\qquad u\circ v\simeq\id_Y
\]
\cite[Section 3]{BCM2006}. The second identity says that $Y$ is
\emph{finitely dominated} by $X$. For a pointed homotopy idempotent
with $f_\#=\id_G$, where $G=\pi_1(X)$, the induced maps $u_\#$ and
$v_\#$ identify $\pi_1(Y)$ with $G$. The universal-cover chains of $Y$
are chain homotopy equivalent to a finite complex $P_*$ of finitely
generated projective $\Z G$-modules. The sum
$\sum_i(-1)^i[P_i]\in K_0(\Z G)$ is its projective Euler class. The
Hattori--Stallings trace of this class records the fixed-point class indices of $f$:
it is the Reidemeister trace, with conjugacy classes inverted when
passing from the left-module convention of \cite[Lemma 4.1]{BCM2006}
to ours. The augmentation trace and the identity coefficient of this projective
class give the ordinary and $L^2$-Lefschetz numbers, respectively
\cite[Lemmas 8.1--8.2]{BCM2006}. Berrick, Chatterji and Mislin use these
relations to deduce the following consequences of the integral support
statement and its numerical trace consequence.

\begin{corollary}[Berrick--Chatterji--Mislin consequences]
\label{cor:homotopy-idempotents}
Let $M$ be a connected closed smooth oriented manifold.
\begin{enumerate}[label=\textup{(\roman*)}]
 \item If $\dim M>2$ and $f:M\to M$ satisfies $f\circ f\simeq f$
 by free homotopy, then some self-map freely homotopic to $f$ has
 exactly one fixed point.
 \item Choose $x_0\in M$ and put $G=\pi_1(M,x_0)$. If
 $f:(M,x_0)\to(M,x_0)$ satisfies $f\circ f\simeq_* f$ by based
 homotopy and $f_\#=\id_G$, then
 \[
  L(f)=L^{(2)}(\widetilde f).
 \]
 Here $\widetilde f$ is the lift to the universal cover fixing a chosen
 point above $x_0$; it is $G$-equivariant. The left side is the ordinary
 Lefschetz number, and the right side is the alternating sum of von
 Neumann traces on the induced $L^2$-homology maps, as in
 \cite[Definition 7.1]{BCM2006}.
\end{enumerate}
\end{corollary}

\begin{proof}
For $x\in K_0(\Z G)$ write $h_C(x)=\HS_{\Z G}(x)(C)$. The
augmentation trace $\epsilon$ is the sum of these coefficients, while
the Kaplansky trace $\kappa$ is the identity coefficient. The sum is
finite, and Corollary~\ref{cor:integral} gives
\[
 \epsilon(x)=\sum_{C\in\con(G)}h_C(x)=h_{[1]_G}(x)=\kappa(x).
\]
This is the weak Bass conjecture for every group. The cited paper uses
projective left modules, whereas we use right modules. The anti-involution
$g\mapsto g^{-1}$, together with transpose on idempotent matrices,
transports between these conventions and relabels the coefficient at
$[g]_G$ by $[g^{-1}]_G$. It preserves identity support and both
$\epsilon$ and $\kappa$. Thus Corollary~\ref{cor:integral} gives the
classical Bass conjecture in the cited convention, and the displayed
equality gives weak Bass there. The forward implications of
\cite[Theorem 1]{BCM2006} and \cite[Theorem 2]{BCM2006} now give
\textup{(i)} and \textup{(ii)}, respectively.
\end{proof}
